Aus $\bar v = \sqrt{8kT/(\pi m)}$ folgt: Die mittleren Schnelligkeiten sind gleich, wenn $T/m$ für beide Gase gleich ist,
\begin{align}
 \frac{T_{\mathrm{O_2}}}{m_{\mathrm{O_2}}} &= \frac{T_{\mathrm{N_2}}}{m_{\mathrm{N_2}}}
\end{align}
Das Verhältnis der Teilchenmassen ist gleich dem Verhältnis der molaren Massen (FoTaBe S.~228):
\begin{align}
 \frac{T_{\mathrm{O_2}}}{T_{\mathrm{N_2}}} &= \TqF = \frac{2\cdot\MOaO}{2\cdot\MNaO} = \Tq \approx \result{\TqP{0}{3}}
\end{align}
Die absolute Temperatur des Sauerstoffs muss also um $\dTrP{0}{2}$ höher sein als die des Stickstoffs (ungefähr im Verhältnis $32:28 = 8:7$).
